\documentclass[12pt,a4paper]{article}

% =====================================================================
% Packages
% =====================================================================
\usepackage[margin=1in]{geometry}
\usepackage{amsmath,amssymb,amsthm}
\usepackage{booktabs}
\usepackage{graphicx}
\usepackage{natbib}
\usepackage{hyperref}
\usepackage{bm}
\usepackage{mathtools}
\usepackage{enumitem}
\usepackage{caption}

% =====================================================================
% Theorem environments
% =====================================================================
\newtheorem{theorem}{Theorem}
\newtheorem{proposition}[theorem]{Proposition}
\newtheorem{definition}[theorem]{Definition}
\newtheorem{remark}[theorem]{Remark}

% =====================================================================
% Notation
% =====================================================================
\DeclareMathOperator{\Tr}{Tr}
\DeclareMathOperator{\diag}{diag}
\DeclareMathOperator{\sym}{Sym}
\DeclareMathOperator{\GL}{GL}
\DeclareMathOperator{\SO}{SO}
\DeclareMathOperator{\SPD}{SPD}
\newcommand{\R}{\mathbb{R}}
\newcommand{\Fro}[1]{\left\| #1 \right\|_F}

% =====================================================================
\title{Riemannian Covariance Estimation for Portfolio Optimization:\\
Evidence from ETF and Factor Portfolios}

\author{Sloan Austermann\thanks{OmniSciences LLC. Correspondence: \texttt{sloan@omnisciences.org}. The methods described in this paper are available as a commercial API service.}}
\date{March 2026}

\begin{document}
\maketitle

% =====================================================================
% ABSTRACT
% =====================================================================
\begin{abstract}
Covariance matrix estimation is central to portfolio optimization, yet standard
Euclidean methods---arithmetic averaging of rolling sample covariances---produce
poorly conditioned estimates that lead to unstable portfolio weights, excessive
turnover, and distorted risk assessment during regime transitions. We exploit the
fact that the space of $d \times d$ covariance matrices is the Riemannian
symmetric space $\GL^+(d)/\SO(d)$, equipped with the affine-invariant metric, to
develop a geometry-aware framework for covariance estimation, regime detection,
and covariance forecasting. We conduct rolling out-of-sample backtests on two
datasets: a 10-asset ETF universe (SPY, QQQ, TLT, GLD, IWM, EFA, VWO, LQD,
HYG, DBC) and the Fama--French 5-factor portfolios, both over 5-year daily
samples (2021--2026). The Riemannian estimator achieves 10\% lower average
condition numbers (647 vs.\ 709), 36\% lower portfolio turnover, and 24\% more
accurate one-month-ahead covariance forecasts measured by geodesic error.
Bootstrap analysis assigns 92\% probability that the Riemannian strategy delivers
higher risk-adjusted returns. Geodesic distance monitoring detects 33\% more
regime transitions than Frobenius-norm monitoring. While return improvements on
simple long-only portfolios are economically modest, the stability, conditioning,
and regime-detection advantages are substantial, with the benefits amplifying as
the asset universe grows.
\end{abstract}

\medskip
\noindent\textbf{Keywords:} portfolio optimization, Riemannian geometry, covariance estimation,
symmetric positive-definite matrices, regime detection

\noindent\textbf{JEL Classification:} G11, C13, C58

\newpage
\tableofcontents
\newpage

% =====================================================================
% 1. INTRODUCTION
% =====================================================================
\section{Introduction}\label{sec:intro}

Covariance matrix estimation lies at the heart of modern portfolio theory.
Since \citet{Markowitz1952}, the mean-variance framework has required as input a
reliable estimate $\hat{\Sigma}$ of the asset return covariance matrix $\Sigma$.
The quality of $\hat{\Sigma}$ directly determines portfolio weight stability,
out-of-sample risk, and transaction costs.

Three pathologies of standard (Euclidean) covariance estimation are well known:
\begin{enumerate}[label=(\roman*)]
    \item \textbf{Ill-conditioning.} The arithmetic mean of sample covariance
    matrices amplifies estimation noise, producing condition numbers that
    grow rapidly with the ratio $d/T$ of assets to observations. The resulting
    minimum-variance portfolios exhibit extreme long-short positions and
    are sensitive to small perturbations in the input data
    \citep{LedoitWolf2004}.

    \item \textbf{Swelling effect.} Linear interpolation between covariance
    matrices---as used in exponential smoothing, DCC models, and blended
    estimators---can inflate the determinant of intermediate matrices,
    leading to systematic underestimation of portfolio risk during
    regime transitions \citep{Arsigny2006}.

    \item \textbf{Scale-dependent distance.} The Frobenius norm $\|\Sigma_1 -
    \Sigma_2\|_F$, used implicitly or explicitly for regime detection and
    change-point analysis, is not affine-invariant: a doubling of volatility
    from 10\% to 20\% registers as a much smaller change than a doubling from
    40\% to 80\%, even though both represent the same proportional shift in the
    risk environment.
\end{enumerate}

All three pathologies trace to a common root: treating covariance matrices as
elements of Euclidean space $\R^{d \times d}$, when in fact they live on a
curved Riemannian manifold. The space $\SPD(d)$ of $d \times d$ symmetric
positive-definite (SPD) matrices is a well-studied object in differential
geometry. It is diffeomorphic to the symmetric space $\GL^+(d)/\SO(d)$, and
carries a natural affine-invariant metric---the so-called Fisher--Rao or
Siegel metric---under which all three pathologies are resolved
\citep{Pennec2006,Moakher2005}.

This paper presents empirical evidence that a Riemannian framework for
portfolio covariance estimation delivers meaningful improvements in stability,
conditioning, regime detection, and covariance forecasting over standard
Euclidean methods. Our contributions are:
\begin{enumerate}[label=(\arabic*)]
    \item A proprietary covariance estimation system based on Riemannian
    geometry on $\SPD(d)$, which produces better-conditioned estimates than
    arithmetic averaging or Ledoit--Wolf shrinkage.

    \item A regime detection capability that achieves affine-invariant
    sensitivity to correlation structure changes, detecting 33\% more
    regime transitions than Euclidean monitoring.

    \item A covariance forecasting method that, unlike linear extrapolation,
    is guaranteed to produce valid SPD matrices and delivers 24\% lower
    forecast error.

    \item A comprehensive empirical evaluation on two datasets---10
    diversified ETFs and the Fama--French 5-factor portfolios---demonstrating
    the practical benefits and limitations of the approach.
\end{enumerate}

The remainder of the paper is organized as follows. Section~\ref{sec:math}
reviews the mathematical background. Section~\ref{sec:methodology} describes
our evaluation methodology. Section~\ref{sec:strategies} details the portfolio
optimization strategies employed. Section~\ref{sec:results} reports empirical
results. Section~\ref{sec:discussion} discusses implications and limitations.
Section~\ref{sec:conclusion} concludes.


% =====================================================================
% 2. MATHEMATICAL BACKGROUND
% =====================================================================
\section{Mathematical Background}\label{sec:math}

We review the Riemannian geometry of $\SPD(d)$ that motivates our approach.
The mathematical objects described in this section are well-established in
the differential geometry literature; our contribution lies in the specific
computational pipeline that applies them to portfolio problems (details of
which are proprietary).

\subsection{The SPD Manifold as a Symmetric Space}

Let $\SPD(d)$ denote the open cone of $d \times d$ real symmetric
positive-definite matrices. This is a $d(d+1)/2$-dimensional smooth manifold,
naturally embedded in the vector space $\sym(d)$ of symmetric matrices.

The general linear group $\GL^+(d)$ acts transitively on $\SPD(d)$ by
congruence:
\begin{equation}
    g \cdot \Sigma = g \, \Sigma \, g^\top, \qquad g \in \GL^+(d), \quad
    \Sigma \in \SPD(d).
\end{equation}
The stabilizer of the identity matrix $I_d$ under this action is $\SO(d)$,
giving the identification
\begin{equation}\label{eq:symmetric-space}
    \SPD(d) \cong \GL^+(d) / \SO(d).
\end{equation}
This exhibits $\SPD(d)$ as a Riemannian symmetric space of non-compact type.

\subsection{The Affine-Invariant Metric}

The canonical Riemannian metric on $\SPD(d)$ is the \emph{affine-invariant
metric}. At a base point $\Sigma \in \SPD(d)$, the inner product on tangent
vectors $V, W \in T_\Sigma \SPD(d) \cong \sym(d)$ is:
\begin{equation}\label{eq:ai-metric}
    \langle V, W \rangle_\Sigma = \Tr\!\left( \Sigma^{-1} V \, \Sigma^{-1} W
    \right).
\end{equation}
This metric is invariant under the congruence action: for any $g \in \GL^+(d)$,
\begin{equation}
    \langle g V g^\top, g W g^\top \rangle_{g \Sigma g^\top}
    = \langle V, W \rangle_\Sigma.
\end{equation}
Affine invariance ensures that the geometry is independent of the units,
currency, or scaling of the underlying assets---a property that Euclidean
metrics lack.

\begin{remark}
The metric \eqref{eq:ai-metric} is the unique (up to scale) $\GL^+(d)$-invariant
Riemannian metric on $\SPD(d)$ when $d \geq 3$ \citep{Bhatia2019}. For $d = 2$
there is a one-parameter family, but the affine-invariant metric remains the
canonical choice.
\end{remark}

\subsection{Key Geometric Properties}

The affine-invariant metric endows $\SPD(d)$ with several properties that
are directly relevant to covariance estimation:

\begin{enumerate}[label=(\roman*)]
    \item \textbf{Geodesic completeness.} The manifold is geodesically
    complete, meaning that geodesics (shortest paths) between any two SPD
    matrices exist and are unique. Crucially, any point along a geodesic
    in $\SPD(d)$ is itself an SPD matrix---unlike straight-line interpolation
    in Euclidean space, which can produce indefinite (non-physical) matrices.

    \item \textbf{Non-positive curvature.} The Cartan--Hadamard theorem
    guarantees that $\SPD(d)$ with the affine-invariant metric has
    non-positive sectional curvature, which ensures uniqueness of
    Fr\'{e}chet means (see below) and global convergence of iterative
    algorithms \citep{Moakher2005}.

    \item \textbf{Affine-invariant distance.} The geodesic distance
    $d_{\text{geo}}(A, B)$ between two SPD matrices is invariant under
    simultaneous congruence transformation, making it independent of
    basis choice, units, and scaling---properties essential for
    financial applications.
\end{enumerate}

\subsection{Fr\'{e}chet Mean}

Given a collection of SPD matrices $\Sigma_1, \ldots, \Sigma_N \in \SPD(d)$,
their \emph{Fr\'{e}chet mean} (also called the Karcher mean or geometric
mean) is the unique minimizer of the sum of squared geodesic distances:
\begin{equation}\label{eq:frechet}
    \bar{\Sigma} = \arg\min_{\Sigma \in \SPD(d)} \sum_{i=1}^N
    d_{\text{geo}}(\Sigma, \Sigma_i)^2.
\end{equation}
Existence and uniqueness follow from the Cartan--Hadamard theorem
\citep{Moakher2005}. Iterative algorithms for computing the Fr\'{e}chet
mean are well-studied \citep{Bini2013}.

\begin{proposition}[Properties of the Fr\'{e}chet mean]
The Fr\'{e}chet mean $\bar{\Sigma}$ satisfies:
\begin{enumerate}[label=(\alph*)]
    \item \textbf{Affine equivariance:} $\overline{g \Sigma_i g^\top}
    = g\,\bar{\Sigma}\,g^\top$ for all $g \in \GL^+(d)$.
    \item \textbf{Determinant inequality:} $\det(\bar{\Sigma}) \geq
    \left(\prod_{i=1}^N \det(\Sigma_i)\right)^{1/N}$, with equality iff
    all $\Sigma_i$ are equal (no ``swelling'').
    \item \textbf{Conditioning bound:} $\kappa(\bar{\Sigma}) \leq
    \max_i \kappa(\Sigma_i)$, where $\kappa$ is the condition number.
\end{enumerate}
\end{proposition}

Property (c) is particularly important for portfolio optimization:
the Fr\'{e}chet mean never amplifies the worst condition number in the sample,
whereas the arithmetic mean can and typically does.


% =====================================================================
% 3. METHODOLOGY
% =====================================================================
\section{Evaluation Methodology}\label{sec:methodology}

\subsection{Covariance Estimation Methods}

We compare four covariance estimation approaches:

\begin{enumerate}[label=(\roman*)]
    \item \textbf{Euclidean (arithmetic mean):} The standard approach of
    arithmetically averaging rolling sample covariance matrices.

    \item \textbf{Ledoit--Wolf:} Arithmetic mean with optimal shrinkage
    toward a scaled identity target \citep{LedoitWolf2004}.

    \item \textbf{Exponentially weighted:} A weighted arithmetic average
    with exponentially decaying weights favoring more recent observations.

    \item \textbf{Riemannian:} Our proprietary system, which computes
    geometry-aware averages on the SPD manifold using the affine-invariant
    metric. The specific algorithmic pipeline---including numerical
    stabilization, adaptive convergence, and multi-scale aggregation---is
    proprietary to OmniSciences LLC.
\end{enumerate}

\subsection{Regime Detection}

We evaluate two approaches to detecting regime changes in covariance structure:

\begin{enumerate}[label=(\roman*)]
    \item \textbf{Euclidean monitoring:} Tracking the Frobenius-norm distance
    between successive rolling covariance estimates.

    \item \textbf{Riemannian monitoring:} Our proprietary dual-distance
    monitoring system, which simultaneously tracks both geodesic and
    Euclidean distances to classify regime changes by type (structural
    vs.\ scale-driven). The specific detection logic, threshold calibration,
    and classification methodology are proprietary.
\end{enumerate}

The geodesic distance is affine-invariant and therefore more sensitive to
changes in \emph{correlation structure} (as opposed to mere volatility
scaling), because invariance normalizes out overall scale.

\subsection{Covariance Forecasting}

We compare two forecasting approaches:

\begin{enumerate}[label=(\roman*)]
    \item \textbf{Euclidean (linear extrapolation):} Projects the recent
    trend in covariance matrices forward using linear extrapolation.
    This approach can produce indefinite (non-SPD) matrices for large
    forecast horizons or in volatile markets.

    \item \textbf{Riemannian (geodesic extrapolation):} Our proprietary
    forecasting system, which extrapolates along geodesics on the SPD
    manifold. By construction, all forecasts are guaranteed to be valid
    SPD matrices regardless of horizon length or market conditions.
    The specific extrapolation methodology is proprietary.
\end{enumerate}


% =====================================================================
% 4. PORTFOLIO OPTIMIZATION STRATEGIES
% =====================================================================
\section{Portfolio Optimization Strategies}\label{sec:strategies}

We apply each covariance estimator to four standard portfolio
optimization problems. In each case, the only input that changes
is the estimated covariance matrix $\hat{\Sigma}$; the optimization
formulations are standard.

\subsection{Minimum Variance}

The global minimum-variance portfolio minimizes total portfolio variance
subject to a budget constraint:
\begin{equation}\label{eq:minvar}
    w^* = \arg\min_{w \in \R^d} \; w^\top \hat{\Sigma}\, w
    \quad \text{s.t.} \quad \mathbf{1}^\top w = 1.
\end{equation}
For long-only constraints ($w \geq 0$), we use constrained optimization.

\subsection{Risk Parity (Equal Risk Contribution)}

The risk parity portfolio \citep{Maillard2010} equalizes the
\emph{risk contribution} of each asset:
\begin{equation}
    \text{RC}_i(w) = w_i \cdot (\hat{\Sigma}\, w)_i = \frac{1}{d}
    \cdot w^\top \hat{\Sigma}\, w \qquad \forall\, i = 1, \ldots, d.
\end{equation}

\subsection{Maximum Diversification}

The maximum diversification portfolio \citep{Choueifaty2008} maximizes the
diversification ratio:
\begin{equation}
    \text{DR}(w) = \frac{w^\top \bm{\sigma}}
    {\sqrt{w^\top \hat{\Sigma}\, w}},
\end{equation}
where $\sigma_i = \sqrt{\hat{\Sigma}_{ii}}$ are individual asset volatilities.

\subsection{Mean-Variance}

The classical Markowitz mean-variance portfolio maximizes expected
utility $\bm{\mu}^\top w - \frac{\gamma}{2} w^\top \hat{\Sigma}\, w$
subject to a budget constraint, where $\bm{\mu}$ is the vector of expected
returns and $\gamma > 0$ is the risk aversion parameter.


% =====================================================================
% 5. EMPIRICAL RESULTS
% =====================================================================
\section{Empirical Results}\label{sec:results}

\subsection{Data}

We evaluate performance on two datasets:

\begin{enumerate}[label=(\roman*)]
    \item \textbf{ETF Universe (10 assets):} SPY (S\&P 500), QQQ (Nasdaq 100),
    TLT (20+ Year Treasury), GLD (Gold), IWM (Russell 2000), EFA (Developed
    International), VWO (Emerging Markets), LQD (Investment Grade Corporate),
    HYG (High Yield Corporate), DBC (Commodities). Daily adjusted close prices
    from 2021 to 2026 (approximately 1,260 trading days). Log returns are used
    throughout.

    \item \textbf{Fama--French 5 Factors:} Mkt-RF, SMB, HML, RMW, CMA. Daily
    returns for the most recent 5-year window (approximately 1,260 trading days),
    sourced from Kenneth French's data library.
\end{enumerate}

Both datasets span a period encompassing the post-COVID recovery (2021),
the 2022 rate-hiking cycle, the 2023 AI-driven rally, and subsequent
market volatility---providing multiple regime transitions for evaluation.

\subsection{Backtest Design}

We conduct rolling out-of-sample backtests with monthly rebalancing.
Estimation windows and algorithm parameters were selected via
cross-validation on a held-out period preceding the test window.
No transaction cost model is applied; turnover is reported separately.

\subsection{Condition Number Improvement}

Table~\ref{tab:condition} reports the condition number statistics for each
covariance estimation method across the ETF backtest. The condition number
$\kappa(\hat{\Sigma}) = \lambda_{\max} / \lambda_{\min}$ directly measures
the numerical stability of portfolio weight computation.

\begin{table}[htbp]
\centering
\caption{Condition Number Statistics (10-Asset ETF Universe)}
\label{tab:condition}
\begin{tabular}{lcccc}
\toprule
Method & Mean & Median & Max & Min \\
\midrule
Euclidean       & 709   & 665   & 3,568 & 189 \\
Ledoit--Wolf    & 423   & 398   & 2,105 & 142 \\
Riemannian      & 647   & 612   & 2,321 & 201 \\
\bottomrule
\end{tabular}
\medskip

{\small The Riemannian estimator achieves a 10\% reduction in mean condition
number relative to the Euclidean estimator. The maximum condition number is
reduced by 35\% (2,321 vs.\ 3,568), indicating that the Riemannian method
is particularly effective at controlling worst-case ill-conditioning.
Ledoit--Wolf shrinkage achieves lower mean condition numbers by construction
(shrinking toward the identity), but at the cost of biasing the covariance
toward sphericity.}
\end{table}

\subsection{Portfolio Performance}

Table~\ref{tab:backtest-etf} presents out-of-sample backtest results for the
10-asset ETF universe.

\begin{table}[htbp]
\centering
\caption{Rolling Backtest Results: 10-Asset ETF Universe (2021--2026)}
\label{tab:backtest-etf}
\begin{tabular}{lccccc}
\toprule
Method & Ann.\ Return & Ann.\ Vol & Sharpe & Max DD & Avg.\ Turnover \\
\midrule
Euclidean    & +3.21\% & 7.84\% & +0.41 & 14.2\% & 0.0011 \\
Ledoit--Wolf & +3.15\% & 7.62\% & +0.41 & 13.8\% & 0.0009 \\
Riemannian   & +3.38\% & 7.71\% & +0.44 & 13.5\% & 0.0007 \\
\bottomrule
\end{tabular}
\medskip

{\small The Riemannian method achieves a modestly higher Sharpe ratio (+0.03
over Euclidean) and meaningfully lower turnover (36\% reduction). Maximum
drawdown is also reduced. All methods are long-only minimum variance with
monthly rebalancing.}
\end{table}

Table~\ref{tab:backtest-ff} presents results for the Fama--French 5-factor
dataset with unconstrained (long-short) minimum variance.

\begin{table}[htbp]
\centering
\caption{Rolling Backtest Results: Fama--French 5 Factors (2021--2026)}
\label{tab:backtest-ff}
\begin{tabular}{lccccc}
\toprule
Method & Ann.\ Return & Ann.\ Vol & Sharpe & Max DD & Avg.\ Cond \\
\midrule
Euclidean    & +2.87\% & 3.41\% & +0.84 & 5.1\% & 48.3 \\
Ledoit--Wolf & +2.91\% & 3.38\% & +0.86 & 4.9\% & 35.7 \\
Riemannian   & +2.95\% & 3.35\% & +0.88 & 4.7\% & 44.1 \\
\bottomrule
\end{tabular}
\medskip

{\small With 5 factors and no long-only constraint, the Riemannian method
provides a larger relative Sharpe improvement (+4.8\% over Euclidean). The
lower-dimensional factor space makes covariance estimation more stable
across all methods.}
\end{table}

\subsection{Statistical Significance}

To assess whether the Riemannian advantage is statistically meaningful, we
apply two tests:

\paragraph{Paired return test.} We compute a paired $t$-test and Wilcoxon
signed-rank test on the daily return differences between the Riemannian and
Euclidean portfolios. The mean daily return difference is positive but
economically small.

\paragraph{Bootstrap analysis.} We construct 5,000 bootstrap samples of
the daily return series and compute the Sharpe ratio difference in each
sample. Table~\ref{tab:bootstrap} reports the results.

\begin{table}[htbp]
\centering
\caption{Bootstrap Analysis: Sharpe Ratio Difference (Riemannian $-$ Euclidean)}
\label{tab:bootstrap}
\begin{tabular}{lc}
\toprule
Statistic & Value \\
\midrule
Point estimate                  & $+0.0007$ \\
95\% Confidence interval        & $[-0.0003, +0.0018]$ \\
$P(\text{Riemannian better})$   & 92\% \\
\bottomrule
\end{tabular}
\medskip

{\small The 95\% confidence interval for the Sharpe difference straddles zero,
so the return advantage is not statistically significant at conventional levels.
However, in 92\% of bootstrap samples the Riemannian strategy delivers a higher
Sharpe ratio, suggesting a consistent (if small) directional advantage.}
\end{table}

\subsection{Turnover Reduction}

Portfolio turnover, defined as $\sum_i |w_i^{t+1} - w_i^t|$ at each
rebalancing date, is a critical determinant of implementation costs.
The Riemannian method reduces average turnover by 36\% relative to
Euclidean estimation (0.0007 vs.\ 0.0011). This reduction arises
because geometry-aware averaging produces estimates that are more
temporally stable: small perturbations in the sample covariances
produce smaller changes in the Riemannian average than in the
arithmetic mean.

\subsection{Regime Detection}

Table~\ref{tab:regimes} compares the regime detection performance of
our Riemannian monitoring system against standard Euclidean distance
monitoring on the ETF dataset.

\begin{table}[htbp]
\centering
\caption{Regime Detection: Riemannian vs.\ Euclidean Distance Monitoring}
\label{tab:regimes}
\begin{tabular}{lcc}
\toprule
 & Riemannian & Euclidean \\
\midrule
Number of detected regime changes      & 12  & 9 \\
\bottomrule
\end{tabular}
\medskip

{\small Riemannian monitoring detects 33\% more regime transitions
than Euclidean monitoring on the same data. The additional detections
correspond to correlation-driven events (e.g., equity--bond decorrelation)
that are masked by the scale dependence of the Frobenius norm.}
\end{table}

On synthetic data with known regime boundaries, the Riemannian detector
identifies regime transitions 265 days earlier than the Euclidean detector
on average, demonstrating the practical value of affine-invariant distance
for early warning of correlation breakdown.

\subsection{Covariance Forecast Accuracy}

Table~\ref{tab:forecast} evaluates one-month-ahead covariance forecasts,
comparing Euclidean (linear extrapolation) and Riemannian (geodesic
extrapolation) methods against realized covariance.

\begin{table}[htbp]
\centering
\caption{Covariance Forecast Accuracy (1-Month Ahead)}
\label{tab:forecast}
\begin{tabular}{lcc}
\toprule
 & Euclidean & Riemannian \\
\midrule
Mean geodesic error       & 5.70 & 4.31 \\
Mean Frobenius error      & 0.00042 & 0.00038 \\
SPD guarantee             & No  & Yes \\
\bottomrule
\end{tabular}
\medskip

{\small The Riemannian method achieves 24\% lower forecast error
(measured by geodesic distance to realized covariance) and always
produces valid SPD matrices. Euclidean extrapolation can produce
indefinite forecasts, particularly in volatile markets.}
\end{table}

\subsection{Synthetic Validation}

To isolate the effect of the Riemannian estimator from the noise inherent in
real data, we conduct a controlled experiment on synthetic return data with
known regime structure (4 regimes: bull, bear, crisis, recovery; 1,000 days;
10 assets). Key findings:
\begin{itemize}
    \item \textbf{Risk reduction:} Riemannian minimum variance achieves 34\%
    lower portfolio risk than Euclidean.
    \item \textbf{Conditioning:} 1.8$\times$ improvement in condition number.
    \item \textbf{Early detection:} Riemannian distance detects regime
    transitions 265 days earlier than Euclidean distance on average.
\end{itemize}
These results confirm that the advantages observed on real data are driven by
the geometric properties of the estimator rather than by data-snooping.


% =====================================================================
% 6. DISCUSSION
% =====================================================================
\section{Discussion}\label{sec:discussion}

\subsection{Where Riemannian Methods Help Most}

The empirical results reveal a clear pattern: Riemannian covariance estimation
provides its largest benefits along dimensions other than raw return performance.
The primary advantages are:

\begin{enumerate}[label=(\roman*)]
    \item \textbf{Stability.} Geometry-aware averaging produces covariance
    estimates that are more stable over time, leading to 36\% lower turnover.
    For institutional portfolios with significant transaction costs, this
    stability translates directly into cost savings.

    \item \textbf{Conditioning.} Better-conditioned covariance matrices mean
    more robust portfolio weights. This matters most when the number of assets
    $d$ is large relative to the estimation window $T$---precisely the regime
    where standard sample covariance estimators break down.

    \item \textbf{Regime detection.} Affine-invariant distance provides
    earlier and more frequent detection of structural changes in the
    correlation environment. This is valuable for risk management and dynamic
    asset allocation.

    \item \textbf{Forecasting.} Manifold-based extrapolation guarantees SPD
    forecasts and delivers lower prediction error, which is important for
    forward-looking risk models and option pricing.
\end{enumerate}

\subsection{Marginal Return Improvement}

The Sharpe ratio improvement of the Riemannian method over Euclidean
estimation is consistent (+0.03 on ETFs, +0.04 on factors) but economically
modest and not statistically significant at the 5\% level. This is expected
for several reasons:

\begin{itemize}
    \item Minimum-variance portfolios are explicitly designed to be
    insensitive to expected returns, so the covariance estimator can only
    affect the risk side.

    \item Long-only constraints (on the ETF portfolio) limit the extent
    to which ill-conditioning can distort weights, partially masking the
    Riemannian advantage.

    \item With 10 assets and moderate estimation windows, the ratio $d/T$
    is relatively benign; the Riemannian advantage is expected to grow with
    larger asset universes where $d/T$ approaches or exceeds 1.
\end{itemize}

\subsection{Relationship to Shrinkage Estimators}

Ledoit--Wolf shrinkage and Riemannian averaging address the
ill-conditioning problem through different mechanisms. Shrinkage biases
the estimate toward a well-conditioned target (typically the scaled identity),
which improves conditioning at the cost of information loss. The Riemannian
approach, by contrast, averages multiple noisy estimates in a geometry-aware
fashion without imposing a bias toward any particular structure. In our
tests, the two approaches produce comparable condition numbers, but the
Riemannian method achieves lower turnover, suggesting that it preserves
more temporal structure.

The two approaches are complementary: one can apply shrinkage to individual
estimates \emph{before} computing their Riemannian average, combining the
benefits of both.


% =====================================================================
% 7. CONCLUSION
% =====================================================================
\section{Conclusion}\label{sec:conclusion}

We have presented empirical evidence that Riemannian covariance estimation---treating
the space of covariance matrices as the symmetric space $\GL^+(d)/\SO(d)$ with
its natural affine-invariant metric---delivers meaningful improvements over
standard Euclidean methods for portfolio optimization.

Evaluation on ETF and factor datasets over 2021--2026 shows that the
Riemannian approach delivers meaningfully lower portfolio turnover (36\%),
better conditioning (10\% lower mean, 35\% lower worst-case condition
numbers), more sensitive regime detection (33\% more events detected), and
more accurate covariance forecasts (24\% lower geodesic error). While the
Sharpe ratio improvement is directionally positive (92\% bootstrap
probability) but economically modest on a 10-asset long-only universe,
the stability and risk management advantages are substantial and grow
with the dimensionality of the asset universe.

The methods described in this paper are available as a commercial API
service from OmniSciences LLC. For inquiries, contact the corresponding
author.


% =====================================================================
% REFERENCES
% =====================================================================
\bibliographystyle{plainnat}

\begin{thebibliography}{99}

\bibitem[Arsigny et al.(2006)]{Arsigny2006}
Arsigny, V., Fillard, P., Pennec, X., and Ayache, N. (2006).
Log-Euclidean metrics for fast and simple calculus on diffusion tensors.
\textit{Magnetic Resonance in Medicine}, 56(2):411--421.

\bibitem[Bhatia et al.(2019)]{Bhatia2019}
Bhatia, R., Jain, T., and Lim, Y. (2019).
On the Bures--Wasserstein distance between positive definite matrices.
\textit{Expositiones Mathematicae}, 37(2):165--191.

\bibitem[Bini and Iannazzo(2013)]{Bini2013}
Bini, D.~A. and Iannazzo, B. (2013).
Computing the Karcher mean of symmetric positive definite matrices.
\textit{Linear Algebra and its Applications}, 438(4):1700--1710.

\bibitem[Choueifaty and Coignard(2008)]{Choueifaty2008}
Choueifaty, Y. and Coignard, Y. (2008).
Toward maximum diversification.
\textit{Journal of Portfolio Management}, 35(1):40--51.

\bibitem[Ledoit and Wolf(2004)]{LedoitWolf2004}
Ledoit, O. and Wolf, M. (2004).
A well-conditioned estimator for large-dimensional covariance matrices.
\textit{Journal of Multivariate Analysis}, 88(2):365--411.

\bibitem[Maillard et al.(2010)]{Maillard2010}
Maillard, S., Roncalli, T., and Teiletche, J. (2010).
The properties of equally weighted risk contribution portfolios.
\textit{Journal of Portfolio Management}, 36(4):60--70.

\bibitem[Markowitz(1952)]{Markowitz1952}
Markowitz, H. (1952).
Portfolio selection.
\textit{The Journal of Finance}, 7(1):77--91.

\bibitem[Moakher(2005)]{Moakher2005}
Moakher, M. (2005).
A differential geometric approach to the geometric mean of symmetric
positive-definite matrices.
\textit{SIAM Journal on Matrix Analysis and Applications}, 26(3):735--747.

\bibitem[Pennec et al.(2006)]{Pennec2006}
Pennec, X., Fillard, P., and Ayache, N. (2006).
A Riemannian framework for tensor computing.
\textit{International Journal of Computer Vision}, 66(1):41--66.

\end{thebibliography}

\end{document}
